\chapter{Esercizi sul livello applicativo}

% ══════════════════════════════════════════════════════════════
\section{Preparare Wireshark}
% ══════════════════════════════════════════════════════════════

Prima di svolgere gli esercizi di laboratorio serve un Wireshark funzionante e utilizzabile \textbf{da utente non-root}
(catturare pacchetti come root è sconsigliato: Wireshark analizza dati che arrivano dalla rete, cioè potenzialmente ostili).

\begin{lstlisting}[language=bash]
sudo apt update
sudo apt install wireshark          # rispondere "Sì" alla cattura per utenti non-root

getent group wireshark              # esiste il gruppo wireshark?
groups                              # l'utente ne fa parte?

sudo dpkg-reconfigure wireshark-common      # se serve cambiare la scelta fatta all'installazione
sudo usermod -aG wireshark $USER            # aggiungere l'utente al gruppo
newgrp wireshark                            # (oppure logout e login)

getcap /usr/bin/dumpcap             # deve mostrare cap_net_raw,cap_net_admin=eip
sudo setcap cap_net_raw,cap_net_admin+eip /usr/bin/dumpcap   # solo se mancano i permessi
\end{lstlisting}

\info{Chi cattura davvero i pacchetti}{
    La cattura non la fa l'interfaccia grafica ma il piccolo programma \texttt{dumpcap}, che ha bisogno delle
    \emph{capabilities} \texttt{cap\_net\_raw} (leggere pacchetti grezzi) e \texttt{cap\_net\_admin} (mettere l'interfaccia in
    modalità promiscua). Assegnarle solo a lui è più sicuro che eseguire Wireshark come root.
}

% ══════════════════════════════════════════════════════════════
\section{HTTP con Wireshark}
% ══════════════════════════════════════════════════════════════

\begin{esercizio}{Interazione base: GET e risposta}
\begin{enumerate}
    \item Avviare il browser e Wireshark, senza iniziare la cattura. Nel campo del filtro di visualizzazione scrivere \texttt{http}.
    \item Attendere poco più di un minuto (il motivo sarà chiaro alla fine), poi avviare la cattura.
    \item Visitare \url{http://www.columbia.edu/~fdc/sample.html} e fermare la cattura.
\end{enumerate}
Si vedranno la richiesta GET e la risposta del server. Espandere solo la sezione HTTP (Frame, Ethernet, IP e TCP chiuse) e
rispondere:
\begin{enumerate}[label=(\alph*)]
    \item Quale versione di HTTP usa il browser? E il server?
    \item Quali lingue il browser dichiara di accettare?
    \item Qual è l'indirizzo IP del proprio computer? E quello del server?
    \item Qual è il codice di stato restituito?
    \item Quando è stata modificata l'ultima volta la pagina?
    \item Quanti byte di contenuto sono stati inviati?
    \item Nei byte grezzi ci sono header che non compaiono nella lista dei pacchetti?
\end{enumerate}
\end{esercizio}

\svolgimento
Una cattura tipica ha questo aspetto (i valori esatti dipendono dal browser e dal momento):

\begin{lstlisting}[language=http]
GET /~fdc/sample.html HTTP/1.1
Host: www.columbia.edu
User-Agent: Mozilla/5.0 (X11; Linux x86_64; rv:128.0) Gecko/20100101 Firefox/128.0
Accept: text/html,application/xhtml+xml,application/xml;q=0.9,*/*;q=0.8
Accept-Language: it-IT,it;q=0.9,en-US;q=0.8,en;q=0.7
Accept-Encoding: gzip, deflate
Connection: keep-alive

HTTP/1.1 200 OK
Date: ...
Server: Apache
Last-Modified: ...
Content-Length: ...
Content-Type: text/html
\end{lstlisting}

\begin{enumerate}[label=(\alph*)]
    \item La versione del browser è in fondo alla \textbf{riga di richiesta}, quella del server all'inizio della \textbf{riga di
          stato}: tipicamente entrambe \texttt{HTTP/1.1}. Non devono coincidere per forza: il server risponde con la versione più
          alta che supporta compatibile con quella del client.
    \item Nell'header \texttt{Accept-Language}. I valori \texttt{q} (\emph{quality factor}, da 0 a 1) indicano la preferenza
          relativa: prima l'italiano, poi l'inglese.
    \item Non sono nel messaggio HTTP: si leggono espandendo la sezione \emph{Internet Protocol}. Nella GET \emph{Source} è il
          proprio host e \emph{Destination} il server; nella risposta sono invertiti.
    \item Nella riga di stato: \texttt{200 OK}.
    \item Nell'header \texttt{Last-Modified} della risposta.
    \item Nell'header \texttt{Content-Length}, che conta i byte del \textbf{corpo}, non dell'intero messaggio.
    \item Sì: la vista compatta mostra solo la prima riga, mentre nei byte grezzi si trovano tutti gli header (\texttt{Date},
          \texttt{Server}, \texttt{Content-Type}, \texttt{User-Agent}, \texttt{Accept-Encoding}, e spesso header aggiunti da proxy
          e cache come \texttt{Via}, \texttt{X-Cache}, \texttt{ETag}).
\end{enumerate}

\info{Perché aspettare un minuto}{
    Il server aggiorna il timestamp \texttt{Last-Modified} ogni minuto, così la pagina sembra modificata da meno di un minuto
    prima del download. È un trucco per \textbf{invalidare le cache} (del browser e dei proxy): un GET condizionale con
    \texttt{If-Modified-Since} riceverà sempre \texttt{200 OK} con la pagina, e mai \texttt{304 Not Modified}.
}

\begin{esercizio}{Un server diverso}
Ripetere l'esperimento con \url{http://scottaaronson.com}. Che differenze si notano nella risposta?
\end{esercizio}

\svolgimento
Le differenze da cercare e commentare sono tipicamente:
\begin{itemize}
    \item \textbf{redirezione}: molti siti rispondono \texttt{301 Moved Permanently} o \texttt{302 Found}, con un header
          \texttt{Location} che punta alla versione \texttt{https://}; si vede quindi una catena di richieste;
    \item un \textbf{software server} diverso (header \texttt{Server}, per esempio \texttt{nginx} invece di \texttt{Apache});
    \item \textbf{assenza di \texttt{Last-Modified}} se la pagina è generata dinamicamente, magari sostituito da un
          \texttt{ETag};
    \item \texttt{Transfer-Encoding: chunked} al posto di \texttt{Content-Length}, quando il server non conosce in anticipo la
          lunghezza della risposta;
    \item molte più GET (fogli di stile, script, immagini), spesso su più connessioni in parallelo.
\end{itemize}

\begin{esercizio}{Una risposta divisa in più segmenti}
Svuotare la cache del browser, avviare la cattura, visitare di nuovo \url{http://www.columbia.edu/~fdc/sample.html} e fermare la
cattura. Il file HTML è abbastanza lungo da non stare in un solo segmento TCP.
\begin{enumerate}[label=(\alph*)]
    \item Quante GET ha inviato il browser? In quale pacchetto si trova la GET della pagina?
    \item In quale pacchetto si trovano il codice di stato e la frase della risposta? Quali sono?
    \item Quanti segmenti TCP con dati sono serviti per trasportare la risposta?
\end{enumerate}
\end{esercizio}

\svolgimento
\begin{enumerate}[label=(\alph*)]
    \item Con la cache vuota c'è una GET per la pagina ed eventualmente altre per gli oggetti collegati (per esempio
          \texttt{favicon.ico}). Il numero di pacchetto è nella prima colonna; il filtro \texttt{http.request.method == "GET"}
          le isola.
    \item È il pacchetto con \texttt{HTTP/1.1 200 OK} nella colonna \emph{Info}. Attenzione: \textbf{non} è il primo pacchetto
          della risposta ma l'\textbf{ultimo}, perché Wireshark mostra il messaggio HTTP solo quando ha riassemblato tutti i
          pezzi; i precedenti sono marcati ``\emph{TCP segment of a reassembled PDU}''.
    \item Selezionando la risposta, la sezione \texttt{[N Reassembled TCP Segments]} elenca tutti i segmenti usati, con i loro
          numeri di pacchetto: il conteggio si legge lì.
\end{enumerate}

\info{Stimare il numero di segmenti a mano}{
    Se il corpo è di $L$ byte, gli header HTTP occupano $H$ byte e la MSS vale 1460 byte, i segmenti di dati sono almeno
    $$\left\lceil \frac{L + H}{\text{MSS}} \right\rceil \qquad \text{per esempio}\quad \left\lceil \frac{30\,000 + 300}{1460} \right\rceil = \lceil 20{,}75 \rceil = 21$$
    Il valore osservato può essere maggiore: il server non riempie sempre i segmenti.
}

% ══════════════════════════════════════════════════════════════
\section{Leggere i messaggi HTTP}
% ══════════════════════════════════════════════════════════════

\begin{esercizio}{Una richiesta e la sua risposta}
Un browser invia la richiesta seguente e riceve la risposta riportata sotto.
\begin{lstlisting}[language=http]
GET /corsi/reti/index.html HTTP/1.1
Host: www.docenti.esempio.it
User-Agent: Mozilla/5.0 (X11; Linux x86_64; rv:128.0) Gecko/20100101 Firefox/128.0
Accept: text/html,application/xhtml+xml
Accept-Language: it-IT,it;q=0.8,en;q=0.5
Connection: keep-alive
If-Modified-Since: Mon, 14 Sep 2026 08:12:00 GMT

HTTP/1.1 304 Not Modified
Date: Tue, 15 Sep 2026 10:02:11 GMT
Server: Apache
Connection: Keep-Alive
Keep-Alive: timeout=5, max=100
ETag: "3f80f-1b6-5f4e2c1a"
\end{lstlisting}
\begin{enumerate}[label=(\alph*)]
    \item Qual è l'URL completo dell'oggetto richiesto?
    \item Il browser chiede una connessione persistente o non persistente?
    \item Si può ricavare dal messaggio l'indirizzo IP del client?
    \item Perché la richiesta contiene \texttt{If-Modified-Since}? Quali risposte sono possibili?
    \item L'oggetto è stato trasferito? Che cosa mostrerà il browser?
    \item Per quanto tempo il server terrà aperta la connessione?
\end{enumerate}
\end{esercizio}

\svolgimento
\begin{enumerate}[label=(\alph*)]
    \item Si ottiene unendo \texttt{Host} e il percorso della riga di richiesta: \url{http://www.docenti.esempio.it/corsi/reti/index.html}.
    \item Persistente: \texttt{Connection: keep-alive} (in HTTP/1.1 la connessione è comunque persistente di default).
    \item No: l'indirizzo IP non fa parte del messaggio HTTP ma dell'header del datagramma IP che lo trasporta. Il messaggio
          applicativo non sa (e non deve sapere) nulla degli strati inferiori.
    \item Il browser (o una cache) possiede già una copia dell'oggetto ricevuta il 14 settembre, e chiede di riceverlo
          \textbf{solo se} è cambiato da allora: è un \textbf{GET condizionale}. Il server risponde \texttt{200 OK} con il nuovo
          oggetto se è stato modificato, \texttt{304 Not Modified} senza corpo altrimenti.
    \item No: la risposta 304 non ha corpo. Il browser mostra la copia che ha in cache, risparmiando banda e tempo.
    \item \texttt{Keep-Alive: timeout=5, max=100}: la connessione viene chiusa dopo 5 secondi di inattività o dopo 100 richieste.
\end{enumerate}

\begin{esercizio}{Cookie passo per passo}
Un utente visita per la prima volta il sito di un negozio online, poi ci torna il giorno dopo dallo stesso browser.
Descrivere gli header relativi ai cookie in ciascun messaggio e che cosa fa il server con il valore del cookie.
\end{esercizio}

\svolgimento
\begin{enumerate}
    \item \textbf{Prima richiesta}: il browser non ha cookie per quel sito, quindi la richiesta non contiene l'header
          \texttt{Cookie}.
    \item \textbf{Prima risposta}: il server crea un identificativo (per esempio \texttt{1678}), crea una riga per quell'utente
          nel proprio \textbf{database di back-end} e risponde con \texttt{Set-Cookie: id=1678}.
    \item Il browser memorizza la coppia (sito, \texttt{id=1678}) nel proprio file dei cookie.
    \item \textbf{Richieste successive} (anche il giorno dopo): il browser allega \texttt{Cookie: id=1678}. Il server consulta il
          database e ``riconosce'' l'utente: carrello, preferenze, raccomandazioni.
\end{enumerate}
HTTP resta \textbf{senza stato}: lo stato è mantenuto agli estremi (file dei cookie nel browser, database nel server) e il
cookie è solo la chiave che li collega.

% ══════════════════════════════════════════════════════════════
\section{Tempi di risposta}
% ══════════════════════════════════════════════════════════════

\begin{esercizio}{DNS, connessioni e pipelining}
Un utente clicca un link. L'indirizzo IP del server non è in cache: per risolverlo si interrogano in sequenza 3 server DNS, con
RTT di 20, 30 e 50 ms. La pagina è composta da un file HTML base e da 8 piccoli oggetti sullo stesso server, con
$\text{RTT}_0 = 40$ ms. I tempi di trasmissione sono trascurabili. Calcolare il tempo per ottenere:
\begin{enumerate}[label=(\alph*)]
    \item il solo file HTML;
    \item la pagina completa con connessioni non persistenti, in sequenza;
    \item la pagina completa con connessioni non persistenti e fino a 5 connessioni parallele;
    \item la pagina completa con una connessione persistente senza pipelining;
    \item la pagina completa con una connessione persistente con pipelining.
\end{enumerate}
\end{esercizio}

\svolgimento
Il DNS costa $20 + 30 + 50 = 100$ ms. Ogni connessione TCP nuova costa 1 RTT di apertura più 1 RTT per richiesta e risposta.
\begin{enumerate}[label=(\alph*)]
    \item $100 + 2\,\text{RTT}_0 = 100 + 80 = \mathbf{180}$ \textbf{ms}.
    \item Ogni oggetto richiede una nuova connessione: $180 + 8 \cdot 2\,\text{RTT}_0 = 180 + 640 = \mathbf{820}$ \textbf{ms}.
    \item Gli 8 oggetti si scaricano in due ``giri'' (5 connessioni, poi 3), ciascuno da $2\,\text{RTT}_0$:
          $180 + 2 \cdot 80 = \mathbf{340}$ \textbf{ms}.
    \item La connessione resta aperta; ogni oggetto costa 1 RTT: $180 + 8 \cdot 40 = \mathbf{500}$ \textbf{ms}.
    \item Le 8 richieste partono una dietro l'altra e le risposte arrivano insieme: $180 + 40 = \mathbf{220}$ \textbf{ms}.
\end{enumerate}

\begin{figure}[H]
\centering
\begin{tikzpicture}
\begin{axis}[xbar, width=10cm, height=5.2cm, xmin=0, xmax=950, bar width=9pt,
    symbolic y coords={{(e) persistente, pipelining},{(d) persistente},{(c) 5 parallele},{(b) non persistente}},
    ytick=data, y tick label style={font=\footnotesize\sffamily}, x tick label style={font=\scriptsize},
    xlabel={tempo (ms)}, label style={font=\footnotesize\sffamily},
    nodes near coords, nodes near coords style={font=\scriptsize\sffamily}, enlarge y limits=0.2, grid=major, grid style={gray!20}]
  \addplot[fill=netblue!55, draw=netblue] coordinates {(220,{(e) persistente, pipelining}) (500,{(d) persistente}) (340,{(c) 5 parallele}) (820,{(b) non persistente})};
\end{axis}
\end{tikzpicture}
\caption{Tempo per scaricare la pagina completa nei quattro casi (DNS incluso).}
\end{figure}

\begin{esercizio}{Distribuire un file: client-server o P2P?}
Un file di $F = 15$ Gbit deve essere distribuito a $N$ utenti. Il server ha capacità di upload $u_s = 30$ Mb/s; ogni utente
scarica a $d_i = 2$ Mb/s e, nel caso P2P, carica a $u_i = 700$ kb/s. Calcolare il tempo minimo di distribuzione con
architettura client-server e P2P, per $N = 10$, $100$, $1000$.
\end{esercizio}

\svolgimento
\textbf{Client-server.} Il server deve caricare $N$ copie ($N F$ bit, al massimo a $u_s$), e ogni client non può scaricare più
velocemente di $d_{\min}$:
$$D_{\text{cs}} \ge \max\left\{ \frac{N F}{u_s},\ \frac{F}{d_{\min}} \right\}$$

\textbf{P2P.} Il server deve caricare almeno \emph{una} copia; il client più lento scarica a $d_{\min}$; e in tutto vanno caricati
$N F$ bit, usando la capacità del server \emph{più quella di tutti i peer}:
$$D_{\text{P2P}} \ge \max\left\{ \frac{F}{u_s},\ \frac{F}{d_{\min}},\ \frac{N F}{u_s + \sum_i u_i} \right\}$$

Con i dati: $F/u_s = 15\cdot 10^9 / 3\cdot 10^7 = 500$ s; $F/d_{\min} = 15\cdot 10^9 / 2 \cdot 10^6 = 7500$ s.

\begin{table}[H]
\centering
\begin{tabular}{rrrr}
\toprule
$N$ & $N F / u_s$ & $N F / (u_s + N u)$ & \\
\midrule
10 & 5\,000 s & $150\cdot 10^9 / 37\cdot 10^6 \approx 4\,054$ s & $D_{\text{cs}} = 7\,500$ s, \ $D_{\text{P2P}} = 7\,500$ s \\
100 & 50\,000 s & $1{,}5\cdot 10^{12} / 10^8 = 15\,000$ s & $D_{\text{cs}} = 50\,000$ s, \ $D_{\text{P2P}} = 15\,000$ s \\
1000 & 500\,000 s & $1{,}5\cdot 10^{13} / 7{,}3\cdot 10^8 \approx 20\,548$ s & $D_{\text{cs}} = 500\,000$ s, \ $D_{\text{P2P}} \approx 20\,548$ s \\
\bottomrule
\end{tabular}
\caption{Tempi minimi di distribuzione.}
\end{table}

Nel client-server il tempo cresce \textbf{linearmente} con $N$; nel P2P ogni nuovo peer porta anche capacità di upload, e il
tempo resta limitato: per $N \to \infty$ tende a $F/u = 15 \cdot 10^9 / 7 \cdot 10^5 \approx 21\,429$ s. È la
\textbf{scalabilità intrinseca} del P2P (la curva vista nel capitolo sulla posta e il P2P).

% ══════════════════════════════════════════════════════════════
\section{DNS con nslookup}
% ══════════════════════════════════════════════════════════════

\begin{esercizio}{Esplorare un dominio}
Eseguire in sequenza, su un dominio a scelta (a lezione \texttt{vatican.va}): una query A, una query NS, una query SOA, una
query A per il name server principale e infine la query A rivolta \textbf{direttamente} al name server principale. Commentare
ciascuna risposta.
\end{esercizio}

\svolgimento
Gli output completi per \texttt{vatican.va} sono riportati nel capitolo sul DNS; qui li riassumiamo.

\begin{table}[H]
\centering
\small
\begin{tabular}{p{4.6cm}p{4.6cm}p{5.4cm}}
\toprule
\textbf{Comando} & \textbf{Risultato} & \textbf{Commento} \\
\midrule
\texttt{nslookup vatican.va} & \texttt{185.152.70.106} & record A. \emph{Non-authoritative}: risponde il resolver locale (\texttt{127.0.0.53}), non un server autorevole \\
\texttt{nslookup -type=NS vatican.va} & \texttt{a.nic.va}, \texttt{b.nic.va}, \texttt{c.nic.va}, \texttt{osiris.namex.it}, \texttt{seth.namex.it} & i server autorevoli (per nome, non per indirizzo) \\
\texttt{nslookup -type=SOA vatican.va} & origin \texttt{a.nic.va}, serial, refresh, \dots & \texttt{a.nic.va} è il primario; nessun glue record in coda \\
\texttt{nslookup -type=A a.nic.va} & \texttt{212.77.0.110} & l'indirizzo del primario \\
\texttt{nslookup -type=A vatican.va 212.77.0.110} & \texttt{185.152.70.106} & chiesto al server autorevole: dovrebbe essere \emph{authoritative}, ma il flag AA può non essere impostato per configurazione \\
\bottomrule
\end{tabular}
\caption{Le cinque query su \texttt{vatican.va}.}
\end{table}

\info{I campi del record SOA}{
    \begin{itemize}
      \item \texttt{serial} (2025062401): versione della zona, incrementata a ogni modifica; i secondari la confrontano per capire
            se aggiornarsi.
      \item \texttt{refresh} (14400 s = 4 ore): ogni quanto un secondario controlla il primario.
      \item \texttt{retry} (3600 s = 1 ora): dopo quanto riprovare se il controllo fallisce.
      \item \texttt{expire} (1209600 s = 14 giorni): dopo quanto un secondario che non raggiunge il primario smette di
            considerarsi autorevole.
      \item \texttt{minimum} (3600 s): TTL minimo, usato anche per le risposte negative (``il nome non esiste'').
      \item \texttt{mail addr} (\texttt{dnsop.net.va}): l'indirizzo dell'amministratore, con il primo punto al posto della @.
    \end{itemize}
}

\begin{esempio}{Il mail server di un dominio}
\begin{lstlisting}[language=bash]
$ nslookup -type=MX mit.edu
mit.edu   mail exchanger = 100 mit-edu.mail.protection.outlook.com.

$ nslookup -type=A mit-edu.mail.protection.outlook.com
Address: 52.101.11.3
Address: 52.101.8.44
Address: 52.101.42.14
Address: 52.101.40.1
\end{lstlisting}
Prima si ottiene il \textbf{nome} del mail server (record MX), poi il suo \textbf{indirizzo} (record A). Il 100 è la
\textbf{preferenza}: con più record MX si prova prima quello con il valore più basso. I quattro indirizzi sono un esempio di
\textbf{distribuzione del carico}: il DNS li restituisce ruotandone l'ordine.
\end{esempio}

\begin{esempio}{Alias e nome canonico}
\begin{lstlisting}[language=bash]
$ nslookup www.mit.edu
www.mit.edu               canonical name = www.mit.edu.edgekey.net.
www.mit.edu.edgekey.net   canonical name = e9566.dscb.akamaiedge.net.
Name:    e9566.dscb.akamaiedge.net
Address: 23.50.153.111
\end{lstlisting}
Una catena di CNAME: \texttt{www.mit.edu} è un alias, e il nome canonico finale è un nodo della CDN Akamai. È il servizio di
\textbf{host aliasing} del DNS, usato dalle CDN per indirizzare ogni utente verso un server vicino.
\end{esempio}

% ══════════════════════════════════════════════════════════════
\section{DNS con Wireshark}
% ══════════════════════════════════════════════════════════════

\begin{esercizio}{Una query standard}
Avviare la cattura, eseguire \texttt{nslookup <nome\_a\_piacere>}, fermare la cattura e filtrare con \texttt{dns}.
\begin{enumerate}[label=(\alph*)]
    \item Qual è la porta di destinazione della richiesta? E la porta sorgente della risposta?
    \item A quale indirizzo IP è inviata la richiesta? È il server DNS locale predefinito?
    \item Di che tipo è la query? La richiesta contiene già delle risposte?
    \item Quante domande e quante risposte contiene il messaggio di risposta?
\end{enumerate}
\end{esercizio}

\svolgimento
\begin{enumerate}[label=(\alph*)]
    \item Entrambe \textbf{53}, la porta ben nota del DNS (sezione \emph{User Datagram Protocol}). La porta di destinazione della
          risposta è invece la porta effimera scelta dal sistema operativo per il client.
    \item Tipicamente \texttt{127.0.0.53} (il resolver locale \texttt{systemd-resolved}), oppure il router di casa o il server
          dell'ISP. Si verifica con \texttt{cat /etc/resolv.conf} o \texttt{resolvectl status}: deve coincidere.
    \item Sezione \emph{Queries}, campo \emph{Type}: \texttt{A (Host Address)}. La richiesta ha \emph{Answer RRs} $= 0$: una
          domanda non contiene risposte.
    \item \emph{Questions} $= 1$ (la domanda viene ripetuta nella risposta, così il client può appaiarle) e \emph{Answer RRs}
          pari al numero di record restituiti, spesso più di uno (CNAME intermedi, più indirizzi).
\end{enumerate}

\info{I campi da guardare nell'header DNS}{
    \begin{itemize}
      \item \textbf{Transaction ID}: appaia richiesta e risposta (ed è ciò che un attaccante dovrebbe indovinare per falsificare
            una risposta).
      \item \textbf{QR}: 0 per le query, 1 per le risposte.
      \item \textbf{RD} (\emph{Recursion Desired}): il client chiede la risoluzione ricorsiva; \textbf{RA} (\emph{Recursion
            Available}): il server la offre.
      \item \textbf{AA} (\emph{Authoritative Answer}): la risposta viene da un server autorevole per il nome.
    \end{itemize}
}

\begin{esercizio}{Una query NS}
Ripetere con \texttt{nslookup -type=NS <nome\_a\_piacere>}.
\begin{enumerate}[label=(\alph*)]
    \item A quale indirizzo è inviata la richiesta? Quante domande contiene? Contiene risposte?
    \item Nella risposta: quante risposte ci sono e che informazioni contengono? Quanti record aggiuntivi, e di che tipo?
\end{enumerate}
Suggerimento: interrogare anche server in alto nella gerarchia (inclusi i root server) con l'opzione \texttt{-norecurse}, e poi
senza, confrontando i risultati.
\end{esercizio}

\svolgimento
\begin{enumerate}[label=(\alph*)]
    \item Lo stesso resolver locale: il tipo di query non cambia il destinatario. Una domanda, di tipo NS, e nessuna risposta.
    \item Tante risposte quanti sono i name server autorevoli del dominio: ognuna contiene il \textbf{nome} di un server. Nella
          sezione \emph{Additional records} ci sono i \textbf{glue record}: record A (e AAAA) con gli \textbf{indirizzi} di quei
          server.
\end{enumerate}

\warning{Perché servono i glue record}{
    Per risolvere \texttt{esempio.it} devo chiedere a \texttt{ns1.esempio.it}; ma per contattare \texttt{ns1.esempio.it} devo
    prima risolvere il suo nome, cosa che richiede di chiedere a \texttt{ns1.esempio.it}\dots{} Il glue record spezza questa
    \textbf{dipendenza circolare} fornendo direttamente l'indirizzo.
}

\info{L'effetto di \texttt{-norecurse}}{
    Con \texttt{-norecurse} il flag RD è 0: il server risponde solo con ciò che sa già. Un root server risponde con un
    \textbf{referral} (i name server del TLD, con i glue record), non con l'indirizzo finale. Senza l'opzione il risultato non
    cambia, perché i root server \textbf{non offrono} la ricorsione (RA $= 0$). Per esempio
    \texttt{nslookup iana.org c.root-servers.net} non restituisce record A, ma i name server del dominio \texttt{org}. (Se si è
    dietro una VPN il risultato può sorprendere: si veda il caso di studio nel capitolo su ICMP.)
}

% ══════════════════════════════════════════════════════════════
\section{DNS: domande d'esame}
% ══════════════════════════════════════════════════════════════

\begin{esercizio}{Cambiare l'indirizzo di un server}
Si cambia l'indirizzo IP di una macchina di nome \texttt{server-1.ortofrutta.it}, che ospita (per esempio) un server FTP.
Spiegare come è gestito l'aggiornamento affinché la modifica sia riflessa nel DNS.
\end{esercizio}

\svolgimento
Che la macchina ospiti un server FTP è irrilevante. L'associazione nome--indirizzo è nel server \textbf{autorevole}, dove
l'amministratore la modifica, ma è \textbf{replicata} nelle cache di moltissimi server. Il DNS non avvisa le cache: il punto
cruciale è il \textbf{TTL}. Nessun aggiornamento attivo; si lascia che i record in cache \textbf{scadano}, e alla richiesta
successiva vengono sostituiti da quelli aggiornati. Nel frattempo alcuni client riceveranno il vecchio indirizzo: errori
temporanei che vengono tollerati.

\info{La strategia pratica}{
    Chi sa di dover cambiare indirizzo \textbf{abbassa il TTL} (per esempio a 60 s) qualche giorno prima, fa la modifica, e
    rialza il TTL a transizione completata. La finestra di incoerenza passa da ore a un minuto.
}

\begin{esercizio}{Risolvere un nome con una cache parziale}
Si deve risolvere \texttt{feijoada.dsc.utfpr.edu.br} da un portatile in un laboratorio di UniNA, con la cache vuota. Il portatile
interroga \texttt{dscna2.unina.it}, nella cui cache (per interrogazioni precedenti) c'è un record NS con l'indirizzo del server
DNS di \texttt{dsc.utfpr.edu.br}. Spiegare come il name server di UniNA gestisce la richiesta e come ci si aspetta che
risponda il server di \texttt{dsc.utfpr.edu.br}.
\end{esercizio}

\svolgimento
\begin{figure}[H]
\centering
\begin{tikzpicture}[yscale=0.9]
  \timeline{p}{0}{portatile}{4.8}
  \timeline{u}{6.2}{\texttt{dscna2.unina.it}}{4.8}
  \timeline{d}{12.4}{server di \texttt{dsc.utfpr.edu.br}}{4.8}
  \draw[msg] (0,-0.5) -- node[lbl, above, sloped] {\circled{1} query ricorsiva: \texttt{feijoada.dsc.utfpr.edu.br}?} (6.2,-1.0);
  \node[lbl, fill=llink!60, rounded corners, anchor=west] at (6.35,-1.6) {\circled{2} in cache: NS di \texttt{dsc.utfpr.edu.br}};
  \draw[msg] (6.2,-2.2) -- node[lbl, above, sloped] {\circled{3} query iterativa} (12.4,-2.7);
  \draw[msg, netgreen!60!black] (12.4,-3.1) -- node[lbl, above, sloped] {\circled{4} risposta autorevole (AA), record A} (6.2,-3.6);
  \draw[msg, netgreen!60!black] (6.2,-3.9) -- node[lbl, above, sloped] {\circled{5} indirizzo IP (messo in cache)} (0,-4.4);
\end{tikzpicture}
\caption{La cache permette di saltare root server, \texttt{br}, \texttt{edu.br} e \texttt{utfpr.edu.br}.}
\end{figure}
\begin{enumerate}
    \item Il portatile invia una query \textbf{ricorsiva}: vuole l'indirizzo finale (o un errore).
    \item \texttt{dscna2.unina.it} consulta la cache e trova il record NS (con l'indirizzo) del server della zona
          \texttt{dsc.utfpr.edu.br}: \textbf{non deve partire dai root server}, e salta i passi root $\rightarrow$ \texttt{br}
          $\rightarrow$ \texttt{edu.br} $\rightarrow$ \texttt{utfpr.edu.br}.
    \item Invia una query \textbf{iterativa} direttamente a quel server, chiedendo il record A di \texttt{feijoada.dsc.utfpr.edu.br}.
    \item Il server è autorevole per la zona: risponde con una \textbf{risposta autorevole} (flag AA) contenente il record A.
    \item \texttt{dscna2.unina.it} mette la risposta in cache per la durata del TTL e la restituisce al portatile.
\end{enumerate}
Il senso della domanda è il ruolo della cache: un solo record NS risparmia 3--4 scambi. È anche il motivo per cui i root
server non collassano: la maggior parte delle risoluzioni non li raggiunge mai.

\begin{esercizio}{Contare i messaggi DNS}
Un host con cache vuota vuole risolvere \texttt{www.cs.esempio.edu}; anche il server DNS locale ha la cache vuota, ed è
autorevole per \texttt{cs.esempio.edu} un server distinto da quello di \texttt{esempio.edu}. Quanti messaggi DNS vengono scambiati
in tutto, se il server locale procede in modo \textbf{iterativo}? Quanti ne gestisce il root server se invece tutta la
catena procede in modo \textbf{ricorsivo}?
\end{esercizio}

\svolgimento
In modo iterativo: host $\rightarrow$ locale (1), locale $\leftrightarrow$ root (2), locale $\leftrightarrow$ TLD \texttt{edu} (2),
locale $\leftrightarrow$ server di \texttt{esempio.edu} (2), locale $\leftrightarrow$ server di \texttt{cs.esempio.edu} (2),
locale $\rightarrow$ host (1): \textbf{10 messaggi}. Il root server gestisce una domanda e una risposta.

In modo ricorsivo i messaggi sono sempre 10, ma ogni server deve \textbf{restare in attesa} della risposta del successivo,
mantenendo lo stato della richiesta: il root server resterebbe impegnato per tutta la durata della risoluzione. È per questo
che i server in alto nella gerarchia rispondono solo in modo iterativo.

\begin{esercizio}{Scrivere i resource record}
La piccola azienda \texttt{ortofrutta.it} ha: un name server \texttt{ns1.ortofrutta.it} (\texttt{203.0.113.2}); un server
\texttt{server-1.ortofrutta.it} (\texttt{203.0.113.10}) che ospita il sito, raggiungibile come \texttt{www.ortofrutta.it}; un
mail server \texttt{mail.ortofrutta.it} (\texttt{203.0.113.25}) che riceve la posta \texttt{@ortofrutta.it}. Scrivere i resource
record nel server autorevole dell'azienda e quelli da inserire nel server del TLD \texttt{it}.
\end{esercizio}

\svolgimento
Nel formato (NAME, VALUE, TYPE), trascurando il TTL:
\begin{lstlisting}
; server autorevole di ortofrutta.it
(ortofrutta.it,            ns1.ortofrutta.it,        NS)
(ns1.ortofrutta.it,        203.0.113.2,              A)
(server-1.ortofrutta.it,   203.0.113.10,             A)
(www.ortofrutta.it,        server-1.ortofrutta.it,   CNAME)
(ortofrutta.it,            mail.ortofrutta.it,       MX)    ; con preferenza, es. 10
(mail.ortofrutta.it,       203.0.113.25,             A)

; server del TLD it
(ortofrutta.it,            ns1.ortofrutta.it,        NS)
(ns1.ortofrutta.it,        203.0.113.2,              A)     ; glue record
\end{lstlisting}
Il record A di \texttt{ns1} nel server del TLD è proprio un \textbf{glue record}: senza, nessuno potrebbe contattare il name
server dell'azienda.

% ══════════════════════════════════════════════════════════════
\section{Socket}
% ══════════════════════════════════════════════════════════════

\begin{esercizio}{Chi parte per primo}
Con i programmi client e server UDP e TCP visti nel capitolo sui socket:
\begin{enumerate}[label=(\alph*)]
    \item che cosa succede se si avvia il client TCP \textbf{prima} del server?
    \item e se si avvia il client UDP prima del server?
    \item un server UDP con un solo socket può servire due client? E un server TCP?
    \item che numero di porta usa il client, se non chiama \texttt{bind()}?
\end{enumerate}
\end{esercizio}

\svolgimento
\begin{enumerate}[label=(\alph*)]
    \item \texttt{connect()} fallisce subito: il SYN del client raggiunge un host dove nessun processo è in ascolto su quella
          porta, e il sistema operativo risponde con un segmento \textbf{RST}. Il programma riceve l'errore
          \emph{Connection refused}.
    \item \texttt{sendto()} ha successo (UDP non verifica nulla), ma il datagramma viene scartato a destinazione (al più torna
          un messaggio ICMP \emph{port unreachable}). Il client resta \textbf{bloccato per sempre} in \texttt{recvfrom()}, a
          meno di impostare un timeout: con UDP l'affidabilità è compito dell'applicazione.
    \item UDP: sì, lo stesso socket riceve datagrammi da chiunque, e \texttt{recvfrom()} dice chi è il mittente a cui rispondere.
          TCP: il server usa un \textbf{welcoming socket} per accettare le connessioni e un \textbf{socket dedicato} per ciascun
          client, restituito da \texttt{accept()}.
    \item Una porta \textbf{effimera} assegnata automaticamente dal sistema operativo (su Linux tipicamente tra 32768 e 60999).
\end{enumerate}

% ══════════════════════════════════════════════════════════════
\section{Domande di teoria}
% ══════════════════════════════════════════════════════════════

\begin{esercizio}{HTTP}
\begin{enumerate}[label=(\alph*)]
    \item Spiegare quali sono le parti di un URL e l'uso di ciascuna.
    \item Spiegare brevemente il three-way handshake usato per stabilire una connessione HTTP.
    \item Illustrare che cos'è il pipelining in HTTP.
    \item Su quali protocolli di trasporto si basa HTTP?
\end{enumerate}
Le risposte devono essere complete e possono essere concise.
\end{esercizio}

\svolgimento
\textbf{(a)} La struttura è \texttt{protocol://[userinfo@]host[:port][/path][?query][\#fragment]}:
\begin{itemize}
    \item \textbf{protocol}: come accedere alla risorsa (\texttt{http}, \texttt{https}, \texttt{ftp}\dots); determina la porta di default;
    \item \textbf{userinfo} (facoltativo, \emph{deprecato}): nome utente e password, che finirebbero in chiaro nei log e nella cronologia;
    \item \textbf{host}: nome (da risolvere con il DNS) o indirizzo IP del server;
    \item \textbf{port} (facoltativa): se omessa, 80 per HTTP e 443 per HTTPS;
    \item \textbf{path}: il percorso della risorsa sul server;
    \item \textbf{query}: dopo \texttt{?}, parametri della richiesta (\texttt{chiave=valore\&\dots});
    \item \textbf{fragment}: dopo \texttt{\#}, una parte della risorsa; è gestito \textbf{solo dal client} e non viene mai inviato al server.
\end{itemize}

\textbf{(b)} Il three-way handshake \textbf{non è di HTTP ma di TCP}, su cui HTTP si appoggia:
\begin{enumerate}
    \item il client invia un \textbf{SYN} con un numero di sequenza iniziale casuale \texttt{client\_isn}, senza dati;
    \item il server alloca buffer e variabili e risponde con \textbf{SYN-ACK}: ack $=$ \texttt{client\_isn}$+1$ e il proprio
          \texttt{server\_isn} casuale;
    \item il client risponde con un \textbf{ACK} (ack $=$ \texttt{server\_isn}$+1$), che può già contenere la richiesta HTTP.
\end{enumerate}
L'apertura costa 1 RTT: una richiesta su una connessione nuova costa circa 2 RTT più il tempo di trasmissione.

\textbf{(c)} Con il \textbf{pipelining} il client invia più richieste \textbf{una dietro l'altra}, senza aspettare le risposte;
richiede una connessione persistente e risparmia RTT. Da HTTP/2 richieste e risposte possono anche essere
\textbf{interlacciate} sulla stessa connessione, con priorità, evitando che un oggetto grande blocchi i successivi.

\textbf{(d)} Tradizionalmente \textbf{TCP} (porta 80, o 443 con TLS). Con \textbf{HTTP/3} anche \textbf{UDP}, attraverso
\textbf{QUIC}, che reimplementa affidabilità e controllo di congestione sopra UDP.

\gbox{\iinfo\ Formule da ricordare}{primary!80!black}{
\begin{itemize}
    \item Connessione non persistente: $2\,\text{RTT} + t_{\text{trasm}}$ per oggetto; persistente: 1 RTT per oggetto dopo il primo;
          con pipelining: 1 RTT per tutti.
    \item Distribuzione di file: $D_{\text{cs}} \ge \max\{N F/u_s,\ F/d_{\min}\}$; \
          $D_{\text{P2P}} \ge \max\{F/u_s,\ F/d_{\min},\ N F/(u_s + \sum u_i)\}$.
    \item DNS: porta 53 su UDP; la cache e il TTL governano la propagazione delle modifiche.
\end{itemize}
}
